\subsection{Global structure at the smallest endpoint}
\label{subsec:smallest-endpoint-global-structure}

At logarithmic order two, ampleness makes the local normalization
calculation global.  The resulting cotangent line and foliation do
not require an Albanese hypothesis.

\begin{theorem}[Global structure at the smallest endpoint]
\label{thm:smallest-endpoint-global-structure}
Let $k$ be algebraically closed of characteristic two, let $X$ be a
smooth integral projective surface, and let $C\subset X$ be a smooth
integral ample divisor.  Put $d=C^2$ and $N=\OO_C(C)$.  Let
$M\ne\OO_X$ have exact order two and satisfy $M|_C\simeq\OO_C$.
Suppose that $M$ has a square root $L$, that its logarithmic form
$\omega$ has exact order two along $C$, and that
\[
 2g(C)-2=3d.
\]
The square-root condition is equivalent to
$[M]\in2\operatorname{NS}(X)$.
Let $\tau:W\to X$ be the $\mu_2$-torsor of $M$,
let $n:V\to W$ be its normalization, put $\pi=\tau n$ and
$B=\pi^*C$, and write $\mathfrak c\subset\OO_V$ for the conductor.
Then the following assertions hold.
\begin{enumerate}
\item The integer $d$ is even and
      $\mathcal T=\omega_C\otimes N^{-3}$ has exact order two.
      The surface $V$ is smooth along $B$, and
      \begin{equation}\label{eq:smallest-endpoint-global-conductor}
       \mathfrak c=\OO_V(-B),\qquad
       \omega_V\simeq\pi^*\bigl(\omega_X(-C)\bigr).
      \end{equation}
      In particular, $V$ is Gorenstein.  The morphism $\pi$ is finite
      flat of degree two, and the sequence
      \begin{equation}\label{eq:smallest-endpoint-global-cover-sequence}
       0\longrightarrow\OO_X\longrightarrow\pi_*\OO_V
       \longrightarrow M^{-1}(C)\longrightarrow0
      \end{equation}
      does not split as a sequence of $\OO_X$-modules.  Moreover,
      \[
       \chi(\OO_V)=2\chi(\OO_X)-d/2.
      \]
\item If $s_C$ is the canonical section of $\OO_X(C)$, then
      $\theta=\omega/s_C^2$ is a global section of
      $\Omega_X^1(-2C)$ whose zero scheme $Z_0$ is finite and disjoint
      from $C$.  It gives a saturated exact sequence
      \begin{equation}\label{eq:smallest-endpoint-global-cotangent}
       0\longrightarrow\OO_X(2C)
       \xrightarrow{\ \theta\ }\Omega_X^1
       \longrightarrow I_{Z_0}\otimes\omega_X(-2C)
       \longrightarrow0,
      \end{equation}
      and
      \begin{equation}\label{eq:smallest-endpoint-global-chern}
       \operatorname{length}(Z_0)=c_2(X)\ge0.
      \end{equation}
\item All global one-forms are rationally collinear with $\omega$;
      more precisely,
      \begin{equation}\label{eq:smallest-endpoint-global-forms}
       H^0(X,\Omega_X^1)
       =\theta\,H^0\bigl(X,\OO_X(2C)\bigr).
      \end{equation}
      Their common annihilator is the saturated, $2$-closed
      rank-one foliation
      \begin{equation}\label{eq:smallest-endpoint-global-foliation}
       \mathcal F
       :=\ker\bigl(\omega:T_X\longrightarrow\OO_X\bigr)
       \simeq\omega_X^{-1}(2C).
      \end{equation}
      It satisfies
      \[
       \mathcal F\cdot C=0,\qquad
       \mathcal F|_C\simeq\mathcal T^{-1}
       =\mathcal T\ne\OO_C,\qquad H^0(X,T_X)=0.
      \]
\item Every smooth integral rational curve $E\subset X$ satisfies
      \begin{equation}\label{eq:smallest-endpoint-global-rational-curve}
       E^2\le-2C\cdot E.
      \end{equation}
      In particular, $X$ has no $(-1)$-curve.
\end{enumerate}
\end{theorem}

\begin{proof}
The lifting argument in the proof of
Theorem~\ref{thm:positive-excess-picard} shows that
$[M]\in2\operatorname{NS}(X)$ is equivalent to the existence of $L$:
the discrepancy from a chosen divisor-class root lies in
$\Pic^0(X)$, whose $k$-points are divisible by two.
The line $L$ has exact order four.  Adjunction gives
$K_X\cdot C=2d$, and the displayed genus equality makes $d$ even.
Corollary~\ref{cor:positive-excess-characteristic-two-boundary}
gives the assertion about $\mathcal T$.
The torsor $W$ is integral: a square root of its defining unit in
$k(X)$ would extend as local unit roots by normality and would
trivialize $M$.

\smallskip\noindent
\emph{The local endpoint data.}
We recall the equality consequences needed below from the proof of
that corollary.  In frames of $M$ obtained by squaring frames of $L$,
the finite model near $C$ is
\begin{equation}\label{eq:smallest-endpoint-global-local-model}
 u_i=b_i^2+t_i^2A_i,\qquad
 z_i=b_i+t_i y_i,\qquad y_i^2=A_i,
\end{equation}
where $t_i=0$ defines $C$ and $u_i,b_i$ are units.
Its inverse curve $B_V$ is integral, and
$p_a(B_V)-g(C)=d/2$.
Lemma~\ref{lem:positive-square-root-genus} gives a reduced integral
Cartier curve $B$ on the normalization, with
\[
 B^2=2d,\qquad
 g(\overline B)=g(C),\qquad
 d=\ell+2\Delta,\quad
 \Delta=p_a(B)-g(C),\quad\ell\ge0.
\]
The suitable resolution in
Lemma~\ref{lem:positive-square-resolution-loss} has smooth strict
transform $D$ satisfying $D^2\ge2d-4\Delta\ge0$.
The logarithmic form of the pulled-back root $L$ has a nonzero
leading conormal coefficient of order zero.  A positive $D^2$
would make this a section of a negative-degree line bundle.
Thus $D^2=0$, $\Delta=d/2$, and $\ell=0$.

For the equality in the local normalization and ambient resolution,
write the successive genus losses as
\[
 e=p_a(B_V)-p_a(B),\quad
 u=p_a(B)-p_a(B_0),\quad
 v=p_a(B_0)-g(C).
\]
They are nonnegative and sum to $d/2$.  The resolution calculation
gives
\[
 0=D^2\ge2d-2u-4v=4e+2u.
\]
Hence $e=u=0$.  The ambient exceptional pullback over $B$ has square
zero and is therefore zero, so $V$ is smooth along $B$.
The finite birational map $B\to B_V$ has zero length quotient
and is an isomorphism.  Nakayama applied to the normalization
extension modulo $t_i$ identifies $V$ with
\eqref{eq:smallest-endpoint-global-local-model} near $B$.
These arguments use the square root of $M$, not algebraic
triviality of $M$.

\smallskip\noindent
\emph{The global conductor and the cover algebra.}
The root-genus proof constructs the rank-one reflexive ideal
\[
 J=\mathfrak c\otimes\OO_V(B)\subset\OO_V.
\]
It is the unit ideal at the generic point of $B$.  Both $J$ and
$\OO_V$ have depth two at closed points of the normal surface $V$.
Thus $Q_0=\OO_V/J$ has no zero-dimensional associated component.
It also has no associated component $B$, so a local equation of
$B$ is a nonzerodivisor on $Q_0$.  Consequently
$J|_B\to\OO_B$ is injective with cokernel of length $\ell=0$.
Nakayama gives $Q_0=0$ on a neighborhood of $B$.

The divisor $B$ is ample, being the finite pullback of $C$.
Any remaining support of $Q_0$ would contain a projective curve,
and such a curve meets $B$.  This contradicts the preceding
vanishing.  Hence $J=\OO_V$ globally.
Finite duality in the root-genus proof now gives
\eqref{eq:smallest-endpoint-global-conductor}.
Since $V$ is Cohen--Macaulay and its dualizing sheaf is invertible,
it is Gorenstein.

Off $C$ the conductor is the unit ideal, so $V$ agrees with the
original torsor there.  Near $C$ it is the finite flat model
\eqref{eq:smallest-endpoint-global-local-model}.  Thus $\pi$ is
finite flat everywhere.  Its algebra has an invertible quotient
$Q$ by the constants.  The original nonconstant summand $M^{-1}$
maps to $Q$; in the displayed basis its generator maps to
$t_i[y_i]$.  This map has zero divisor exactly $C$ and is invertible
off $C$, so $Q\simeq M^{-1}(C)$.

If \eqref{eq:smallest-endpoint-global-cover-sequence} split, its
complementary line $Q$ would have linear changes of local generators.
The square of such a generator has zero component in $Q$, because
this is so in the generic purely inseparable degree-two field.
It would therefore define a map
\[
 Q^{\otimes2}\longrightarrow\OO_X.
\]
This map is a section of $Q^{-2}\simeq\OO_X(-2C)$, hence is zero.
A nonzero generic generator would have square zero in a field,
a contradiction.  This proves nonsplittability.  Finally, $M$ is
numerically trivial because it is torsion, and Riemann--Roch gives
\[
 \chi(M^{-1}(C))
 =\chi(\OO_X)+\frac{d-K_X\cdot C}{2}
 =\chi(\OO_X)-d/2.
\]
The Euler-characteristic formula follows from the sequence.

\smallskip\noindent
\emph{The saturated cotangent line.}
Divisibility of a section by a power of the smooth prime divisor
$C$ can be checked at its generic point in the regular local rings
of $X$.  Hence $\theta=\omega/s_C^2$ is global.
Near $C$, the local expression is $dA_i/u_i$.
Smoothness of $y_i^2=A_i$ over every point of $C$ says that $dA_i$
is nonzero there, so $\theta$ has no ambient zero on $C$.
Off $C$, the original hypersurface torsor is normal.
A divisorial zero of $\omega=du_i/u_i$ would make $z_i^2=u_i$
singular along its inverse divisor, contradicting regularity in
codimension one.  Thus $\theta$ has no divisorial zero anywhere,
and its zero scheme $Z_0$ is finite and disjoint from $C$.

If $f,g$ are its coefficients in a local cotangent frame, wedge
with $\theta$ has image the ideal $(f,g)$ times the determinant
line.  The absence of a common height-one factor gives its kernel
as the line generated by $\theta$.  This proves
\eqref{eq:smallest-endpoint-global-cotangent}; its quotient is
torsion-free, so the first line is saturated.
At a zero, $f,g$ form a regular sequence in the regular local
surface ring.  The zero-scheme formula and $K_X\cdot C=2d$ give
\[
 \operatorname{length}(Z_0)
 =c_2\bigl(\Omega_X^1(-2C)\bigr)
 =c_2(X)-2K_X\cdot C+4d=c_2(X).
\]

\smallskip\noindent
\emph{Global forms and vector fields.}
If $\omega\wedge\beta\ne0$ for a global one-form $\beta$, division
by $s_C^2$ gives a nonzero section of $H=\omega_X(-2C)$.
Here $H\cdot C=0$ and $H|_C=\mathcal T$.
Since the nontrivial degree-zero line $\mathcal T$ has no nonzero
section, its zero divisor contains $C$.  Such an effective divisor
has positive intersection with $C$: its $C$-component contributes
$d>0$, and the remaining components have nonnegative intersection.
This contradicts $H\cdot C=0$.
All global forms are therefore rationally collinear with $\omega$.
Their images in the torsion-free quotient of
\eqref{eq:smallest-endpoint-global-cotangent} vanish generically
and hence everywhere, proving
\eqref{eq:smallest-endpoint-global-forms}.

Dualizing the cotangent sequence and using
$\mathcal Hom(I_{Z_0},\OO_X)=\OO_X$ identifies its annihilator
with $\omega_X^{-1}(2C)$.
Multiplication by $s_C^2$ does not change the sheaf kernel, so
this is precisely the kernel in
\eqref{eq:smallest-endpoint-global-foliation}.
The inclusion into $T_X$ is a line subbundle away from $Z_0$,
although it may fail to be a subbundle at $Z_0$.
Locally $\omega=du/u$, and a derivation $D$ lies in its kernel
exactly when $D(u)=0$.  Commutators and restricted squares of such
derivations again annihilate $u$.  This proves the foliation and
$2$-closedness assertions.  Adjunction gives
\[
 \mathcal F\cdot C=0,\qquad
 \mathcal F|_C\simeq\omega_C^{-1}\otimes N^3
 =\mathcal T^{-1}.
\]

For a global vector field $v$, the function $\omega(v)$ is a
constant on the projective integral surface and vanishes on $C$.
Thus $v$ is a section of $\mathcal F$.
A nonzero section of this line would have an effective zero
divisor of intersection zero with the ample curve $C$.
It would consequently be nowhere zero and would trivialize
$\mathcal F$, contrary to its restriction to $C$.
Hence $H^0(X,T_X)=0$.

Finally, for a smooth rational curve $E$, the restriction
$\OO_E(2C)\to\Omega_X^1|_E$ is nonzero because $\theta$ has no
divisorial zero.  Its composite with
$\omega_E=\OO_{\PP^1}(-2)$ is zero by degree.
The conormal sequence gives a nonzero map
$\OO_E(2C)\to\OO_E(-E)$, proving
\eqref{eq:smallest-endpoint-global-rational-curve}.
Since $C\cdot E\ge1$, a $(-1)$-curve is impossible.
\end{proof}

\begin{remark}[Positive self-intersection and global forms]
\label{rem:smallest-endpoint-positive-square-forms}
In the hypotheses of
Theorem~\ref{thm:smallest-endpoint-global-structure}, replace
ampleness of $C$ by $C^2>0$, retaining smoothness of $X$.
Then rational collinearity of all global one-forms and
$H^0(X,T_X)=0$ still hold.
Indeed, the endpoint proof still gives the nontrivial line
$\mathcal T$ and a divided form $\theta$ with no ambient zero
along $C$.  The preceding argument for a section of
$\omega_X(-2C)$ used only $C^2>0$, so proves collinearity.

For a global vector field $v$, the same constant-contraction
argument gives $\omega(v)=0$.  Near $C$, its annihilator is a
line bundle with restriction $\mathcal T^{-1}=\mathcal T$.
If $v\ne0$ and its exact order along $C$ in this local line is $j$,
division by $C^j$ gives a nonzero first coefficient in
\[
 H^0\bigl(C,\mathcal T\otimes N^{-j}\bigr).
\]
For $j=0$ this contradicts nontriviality of $\mathcal T$;
for $j>0$ the degree is negative.
Thus $v=0$.  This argument does not establish the global conductor
or the exact global cotangent sequence without ampleness.
\end{remark}
